% 本文件由 assemble.py 按 _work/plans/ch17.json 逐字组装，请勿手改（改 plan 或 blocks 后重新生成）。
\chapter{原子物理}
\label{ch:17}

\section{本章知识框架}
% ---- block 002-01 (原稿第2页) kind=summary ----
\textbf{原子物理}
\begin{itemize}
\item 光电效应\quad $E_k=h\nu-W_0=\dfrac{hc}{\lambda}-W_0$\quad $h\nu_c=W_0$\quad $eU_c=E_k$\quad $E_k=\dfrac12m_ev^2$
\item 爱因斯坦光电效应方程\quad 光子 $\varepsilon=h\nu$\quad $h=6.63\times10^{-34}\unit{J\cdot s}$\quad $hc=2\times10^{-25}$ % CHECK: 首字形似“受”，按语义转为“爱”；hc 后原稿未写单位
\item 光的波粒二象性
  \begin{itemize}
  \item[$\to$] 干涉、衍射、偏振 % 原稿：该箭头从“光的波”上方引出
  \item[$\to$] 光电效应、康普顿效应、光有动量、能量、 % 原稿：该箭头从“粒”下方引出
  \end{itemize}
\end{itemize}

% ---- block 002-05 (原稿第2页) kind=summary ----
\textbf{原子结构，原子结构} % 页面左侧标注(其右为大括号，含下面“原子结构”“原子核”两支)；CHECK: 第二个“原子结构”或应为“原子核”，原样转录
\begin{itemize}
\item 原子结构\\
  \begin{tabular}[t]{@{}lll@{}}
  汤姆孙发现电子 & 氢原子光谱 & \\
  原子的核式结构 & 玻尔理论 & 氢原子能级、能级公式、 \\
  \end{tabular}
\item 原子核\\
  \begin{tabular}[t]{@{}ll@{}}
  贝克勒尔、天然放射现象、 & 放射性同位素的应用与防护 \\
  原子核的组成 & 核力和核能 \\
  原子核的衰变 & 裂变反应和聚变反应 \\
  \end{tabular}
\end{itemize}


\section{黑体辐射与能量子}
% ---- block 005-01 (原稿第5页) kind=knowledge ----
\textbf{选5} % 页面左上角的小标注

\noindent 热辐射 $\propto\left\{\begin{array}{@{}l@{}}\text{温度}\\ \text{材料}\\ \text{\unclear{表面}状况}\end{array}\right.$\quad （任一物体）辐射电磁波 % CHECK: “∝”手写形似“α”，按语义转为∝；括号内第三项“表面状况”被覆写，字迹不清

\noindent 黑体：吸收所有的入射电磁波而不反射的物体（小孔）\quad 不黑：太阳、窗户\quad 只辐射不反射

\noindent 对于$\langle$绝对黑体$\rangle$热辐射$\langle$只与$T$有关$\rangle$

\noindent $T\uparrow$\ 强度$\uparrow$\quad 峰值向短波移动

\noindent （红的发紫）

\notefig[0.3]{figures/p005_a.png}

% ---- block 005-03 (原稿第5页) kind=knowledge ----
\noindent 普朗克\ $\varepsilon=h\nu$（能量子）\quad $h=6.63\times10^{-34}\unit{J\cdot s}$\quad 不连续$=$量子化

% ---- block 005-05 (原稿第5页) kind=knowledge ----
\noindent 电子 $1e$ 经 $U$ 后\\
\hspace*{2em}$W=-eU$\quad 自带负电

\noindent 电子伏特\quad $1e\cdot1\mathrm{V}=1.6\times10^{-19}\mathrm{C}\cdot1\mathrm{V}=1.6\times10^{-19}\mathrm{J}=1\mathrm{eV}$

\[
\underset{\text{光子能量}}{\text{可见光范围}}\quad \overset{\text{红外}}{1.6\unit{eV}}-\overset{\text{紫外}}{3.1\unit{eV}}
\] % 原稿：“光子能量”写在“可见光范围”下方；“红外”“紫外”两个小字分别写在“1.6eV”“3.1eV”上方(位置判断，可核对)

% ---- block 188-01 (原稿第188页) kind=note ----
\textbf{原子物理 和 天体〉物理}% CHECK: 标题中的“〉”疑为补写“物理”的插入符

1、黑体 $T\uparrow$ $\lambda$ 极大值左移% CHECK: 原稿此处为“〉”形符号，疑为 λ（与右侧 p=h/λ 中的 λ 写法相同）

\notefig[0.25]{figures/p188_a.png}

$p=\dfrac h\lambda$ 散射后 $\lambda\uparrow$ $p\downarrow$．


\section{光电效应}
\subsection{光电效应的规律与方程}
% ---- block 002-02 (原稿第2页) kind=knowledge ----
\textbf{光电效应，波粒二象性} % 页面左侧标注(总标题，对应下面三个专题：光电效应规律、光电效应图像、波粒二象性与物质波)

光电效应规律的理解及应用.
\begin{itemize}
\item[] \unclear{入射}光强度大 $\to$ 光子数目多 $\to$ 发射光电子多 $\to$ 饱和电流大
\item[] 光子频率高 $\to$ 光子能量大 $\to$ 光电子的最大初动能大
\end{itemize}

% ---- block 005-06 (原稿第5页) kind=knowledge ----
\begin{itemize}
\item 赫兹$\langle$实验\unclear{发现}$\rangle$光电效应 % 原稿：“发现”两字与其它笔迹重叠，不清；该行与下一行前有大括号
\item 爱因斯坦：引入光子，解释光电效应（粒子性）\quad $\varepsilon=h\nu=h\dfrac{c}{\lambda}$
\end{itemize}

\noindent 光强/功率\quad $P=n\cdot h\nu$\quad （$n$为$\overset{\text{1s}}{\text{单位时间}}$$\overset{\text{光电子数}}{\text{光子数}}$） % 原稿：“单位时间”上方小注“1s”，“光子数”上方小注“光电子数”

\noindent 逸出功\quad 由金属决定\quad 同一金属 $W_0$ 相同

\noindent 光子（光照能量）$h\nu=h\dfrac{c}{\lambda}$\quad 仅由 $\nu$ 决定\quad $\neq$颜色

\noindent $\langle$最大初动能$\rangle$ $E_k=h\nu-W_0=eU_c$\quad 当 $h\nu=W_0$ 时\quad $\nu=\dfrac{W_0}{h}\to$ 无光电子/光电流

\noindent $\nu_0=\dfrac{W_0}{h}$ $\underset{\text{(极限频率)}}{\text{为截止频率}}$\quad 由金属决定\quad $\nu<\nu_0$ 时 光电流为0\quad 一个光子对应一个光电子

\noindent 光子 $\leftrightarrow$ 电子 一对一\quad
\begin{tabular}[t]{@{}lll@{}}
 & & 强蓝光\ 弱\unclear{蓝}光\\
个体 & 频率---决定初动能 & $\nu$相同\\
整体 & 光强---决定电子数/光子数 & $P$\textsubscript{光强}不同\\
\end{tabular} % 原稿：“个体”“整体”后接大括号；右侧“强蓝光 弱蓝光/ν相同/P(光强)不同”为小注

\noindent 发生条件\quad $h\nu>W_0$ 或 $\nu>\nu_0$

% ---- block 005-09 (原稿第5页) kind=knowledge ----
\noindent \textbf{发生要求}
\begin{itemize}
\item 大于极限频率
\item 小于 最大波长
\end{itemize}

\noindent 同一种光 $U_c$ 相同\\
同一种金属 $W_0$ 相同

% ---- block 008-06 (原稿第8页) kind=knowledge ----
\notefig[0.25]{figures/p008_c.png}

只看最小$\lambda$ % 图 p008_c 内亦含此字

$\lambda<\lambda_0$就可以产生光电子

% ---- block 192-01 (原稿第192页) kind=knowledge ----
1.\ 发生光电效应\ $h\nu>W_0$\quad 跃迁\ $E_{nm}=E_n-E_m$\ 正好\quad 电离\ $E>E_m-E_n$ % 原稿跃迁式中E下方的下标难辨(nm/mn)

$E_k=h\nu-W_0$\hfill 使$R\uparrow$\ $U_{\text{遏止}}\uparrow$ % 右侧“使R↑ U遏止↑”为写在该行右上方的批注
\notefig[0.25]{figures/p192_a.png}

黄光照\textcircled{A}有示数\quad $\nu_{\text{黄光}}>\nu_{\text{极限}}$ % 原稿此处符号写得形似U/ν，按频率ν转录

增加光强\ \textcircled{A}变大\qquad 该装置可测得从阴极板上逸出的光电子个数

若移动$P$电流表示数可能不变（可能已达饱和光电流） % “移动”二字原稿为写在“若”上方的小字

将$P$移至最左侧，阴阳极板间电压为0，阴极逸出的光电子部分能到达阳极，\textcircled{A}不为0。在某个范围内无光电流（达遏止电压时）。

光强大（频率一定时）$\to$光子数目多$\to$发射出电子多$\to$饱和光电流大

光子频率高$\to$光子能量大$\to$光电子最大初动能大。

% ---- block 285-02 (原稿第285页) kind=formula ----
% 页面上缘被裁，此块上方内容看不到
\begin{gather*}
eU_0-W_0=E_K\\
eU_1+h\nu_1-W_0=E_K
\end{gather*}
% CHECK: 第一式“eU_0”也可能是别的写法（字形不清）

% ---- block 189-07 (原稿第189页) kind=note ----
硅吸收光子 不是光电效应．PN 结作用下形成电流——光伏效应% 原稿“不是光电效应”及“PN结…光伏效应”下有下划线

\subsection{光电效应图像}
% ---- block 002-03 (原稿第2页) kind=diagram ----
光电效应图像的理解及应用.

\notefig[0.25]{figures/p002_a.png}
\noindent $k=h$\qquad $W_0=|-E|$

\notefig[0.25]{figures/p002_b.png}

% 下图为I--U曲线(黄光、蓝光两条)；图左上手写的“h=ke”已含在此图中
\notefig[0.3]{figures/p002_c.png}
\noindent 高强光 左大频 % CHECK: “高强光”首字形似“高”，语义上疑为“同”；另I--U图中“黄光”曲线(上)起点比“蓝光”更靠左，与“左大频”似有矛盾，原样转录

% ---- block 005-12 (原稿第5页) kind=diagram ----
\notefig[0.3]{figures/p005_e.png}
\noindent $E_k=h\cdot\nu-W_0$ % 以上图中已含(含纵轴截距标注“-W_0”)

\notefig[0.3]{figures/p005_f.png}
\noindent $U_c=\dfrac{h}{e}\cdot\nu-\dfrac{W_0}{e}$ % 以上图中已含(含纵轴截距标注“-W_0/e”)

\notefig[0.3]{figures/p005_g.png}
\noindent 高强光 左大频 % CHECK: 同第2页，首字形似“高”，语义疑为“同”


\section{康普顿效应、光子动量与光压}
\subsection{康普顿效应与光子动量}
% ---- block 006-01 (原稿第6页) kind=knowledge ----
\textbf{康普顿效应}\quad 波长$\lambda_0$的X射线照射石墨，散射成分中出现波长大于$\lambda_0$的成分。

光子与电子作用时，动量能量都守恒\quad $P=\dfrac{h}{\lambda}$

% ---- block 190-02 (原稿第190页) kind=knowledge ----
2.\ \unclear{康普顿}—粒子性—动量

% ---- block 040-03 (原稿第40页) kind=formula ----
\[
\left\{\begin{aligned}
E&=mc^2\\
E=h\nu&=\frac{hc}{\lambda}\quad\to\ p=\frac Ec=\frac h\lambda\\
p&=mc
\end{aligned}\right.
\]

% ---- block 039-05 (原稿第39页) kind=method ----
正负电子湮灭释放光子
\notefig[0.35]{figures/p039_c.png}
\[
2mc^2=\frac{2hc}{\lambda}\qquad\lambda=\frac{h}{mc}\qquad E\uparrow\ \lambda\downarrow
\]
质能 $2mc^2$\qquad 总能量$=2mc^2+2E_k$

\subsection{光压}
% ---- block 039-04 (原稿第39页) kind=method ----
\textbf{光压强I类}\quad $p=\dfrac{h\nu}{c}=\dfrac h\lambda=\dfrac Ec$\qquad 最大频率 $\dfrac{eU_c}{h}$
\notefig[0.3]{figures/p039_b.png}
吸光\ $\Delta P_1=P$\qquad 反射光\ $\Delta P_2=2P$

白板受冲击力大

% ---- block 040-01 (原稿第40页) kind=derivation ----
$p=\dfrac Ec$

以功率$P_0$射光
\notefig[0.3]{figures/p040_a.png}
光照\ \red{被反射}\quad \red{$\Delta P=2P$} % “被反射”被红笔圈出，红笔写 ΔP=2P
\begin{gather*}
-\Delta F\,\Delta t=-P-P % CHECK: 左端首个符号似“-”，按物理意义(ΔFΔt=2P)照此转录
\\
\Delta F=\frac{2P}{\Delta t} % 此行末尾原稿有一个小记号
\\
\Sigma\Delta F=\Sigma\frac{2E}{c\Delta t}=\frac{2P_0}{c}\qquad\therefore F=\frac{2P_0}{c}
\end{gather*}
光压强 $I=\dfrac{2P_0}{cS}$

% ---- block 040-02 (原稿第40页) kind=derivation ----
功率$P_0$射光\ \red{全被吸收}\quad $\Delta P=P$ % “全被吸收”被红笔圈出，并有黑色波浪线
\notefig[0.25]{figures/p040_b.png}
取$\Delta t$内光子数 $n=\dfrac{P_0\Delta t}{h\,\frac c\lambda}$\qquad $F\Delta t=nP$\quad 故 $I=\dfrac FS=\dfrac{P_0}{cS}$

% ---- block 295-02 (原稿第295页) kind=derivation ----
每秒每平\unclear{米}能量 $E$\quad 光子\unclear{垂直}反弹
\begin{align*}
p&=\frac{2E_0}{c\Delta t}\\ % CHECK: 分母含 Δt，原样转录（光子动量改变量应为 2E_0/c）
ma&=n\frac{\Delta p}{\Delta t}=\frac{ES\Delta t}{E_0}\cdot\frac{2E_0}{c\Delta t}
\end{align*}
加速度\textsuperscript{大小}写时写成正值。

% ---- block 040-04 (原稿第40页) kind=problem ----
% 原稿此处有一条红色斜线横贯页面(分隔其上、其下两部分内容)
\begin{problem}[{①\ 光帆离开太阳系}]
\notefig[0.25]{figures/p040_c.png}
太阳 $G$\unclear{引}质量

（太阳图标）$M$\ 太阳单位时间辐射能量为$P$，光帆与光垂直，将太阳光一半反射一半吸收，求光帆$S$取值满足的条件 % 第一行末尾“吸”字被页边切断，第二行行首又写“吸收”
\begin{solution}
\[
\left\{\begin{aligned}
I&=\frac{-\frac P2-P}{cS}\qquad S=4\pi r^2\ \text{（球面波辐射）}\\
\red{F_{\text{光}}=}\,IS_{\text{帆}}&>\frac{GMm}{r^2}\qquad\textcircled{力}
\end{aligned}\right.
\qquad\therefore\ S_{\text{帆}}>\frac{8\pi cGMm}{3P}
\]
\[
F=P\cdot S
\]
\end{solution}
\end{problem}

% ---- block 040-05 (原稿第40页) kind=problem ----
\begin{problem}[{②}]
激光器功率$P_0$，光束截面积$S$，垂直照射被全部反射时，光对物体压力为$F=2PN$\ ↑\unclear{数}\quad $P$为动量，求光压 % $N$ 下方有一向上箭头，箭头下有一字，似“数”
\begin{solution}
$E=pc$\qquad $P_0=NE$ % 此行 E=pc 处有一条弧形箭头指回左侧
\[
I=\frac FS=\frac{2PN}{S}=\frac{2P_0}{cS}
\]
右侧草算：
\begin{gather*}
E=pc\\
P_0=NE=Npc\\
\boxed{N=\frac{P_0}{pc}} % 原稿 N=P_0/(pc) 被圈出
\end{gather*}
\red{$\checkmark$} % 右侧有一个红笔对勾
\end{solution}
\end{problem}

% ---- block 040-06 (原稿第40页) kind=method ----
\textbf{③}\ 激光动力机.

\red{$F=$}\,$IS=ma$\ \unclear{同}牛二
\begin{itemize}
\item $I=\dfrac{2P_0}{cS}$\quad 光压 % CHECK: 若 P_0 为单位面积辐射功率，则 I=2P_0/c；原稿写 2P_0/(cS)，照录
\item $P=P_0\cdot S$\quad $P_0$为单位面积辐射功率
\item $S$为薄膜面积
\end{itemize}


\section{波粒二象性与物质波}
% ---- block 006-02 (原稿第6页) kind=summary ----
\begin{itemize}
\item 粒子性
  \begin{itemize}
  \item 两个效应
    \begin{itemize}
    \item 光电效应
    \item 康普顿效应
    \end{itemize}
  \item 三个量
    \begin{itemize}
    \item 动量
    \item 能量
    \item 电量
    \end{itemize}
  \end{itemize}
\end{itemize}

\begin{itemize}
\item 粒子性：少，$\lambda$小，与物质作用
\item 波动性：多，$\lambda$大，光传播过程
\end{itemize}

\begin{itemize}
\item 粒子波动性
  \begin{itemize}
  \item 实物波（德布罗意波）
  \item 概率波（微观粒子）
  \end{itemize}
\item 波动性
  \begin{itemize}
  \item $\lambda=\dfrac{h}{p}$\quad $\nu=\dfrac{\varepsilon}{h}$\qquad 2个量
  \item 干涉\quad 衍射\quad 偏振\qquad 3个现象
  \end{itemize}
\end{itemize}

电子动量\quad $P=\sqrt{2mUe}$ % 原稿位于页面右上侧（波动性大括号上方）；首字“电”有涂改叠写

\begin{itemize}
\item 应用
  \begin{itemize}
  \item 电子显微镜（电子$m$大\ $P$大\ $\lambda$小\ 难衍射） % 原稿“难衍射”写在“λ小”下方一行，其后接右括号
  \item 热慢中子衍射 % 原稿“热”与“慢”之间被大括号的竖线隔开
  \item 提高精度：质子显微镜，增大速度
  \end{itemize}
\end{itemize}

% ---- block 002-04 (原稿第2页) kind=knowledge ----
\textbf{对波粒二象性、物质波的理解.}
\begin{itemize}
\item 数量上\quad 少 粒子性\quad 多 波动性
\item 频率上\quad 高 粒子性\quad 低 波动性
\item 传播与作用上\\
  $\uparrow$ 波动性\quad $\uparrow$ 粒子性 % 原稿：两个向上箭头分别在“波动性”“粒子性”上方，指向上一行的“传播与作用上”
\item 波动性和粒子性的统一\quad $\varepsilon=h\nu$\quad $p=\dfrac{h}{\lambda}$
\item 物质波\quad $\lambda=\dfrac{h}{p}=\dfrac{h}{mv}$
\end{itemize}

% ---- block 192-07 (原稿第192页) kind=knowledge ----
光的波动性\fbox{不是}光子相互作用造成的。

% ---- block 189-06 (原稿第189页) kind=formula ----
$E=\dfrac{p^2}{2m}$，\quad $p=\sqrt{2mE}$，\quad $\lambda=\dfrac hp=\dfrac{h}{\sqrt{2mE}}$

% ---- block 008-07 (原稿第8页) kind=knowledge ----
波粒二象性：$P=\dfrac{h}{\lambda}$\qquad $\lambda=\dfrac{h}{p}=\dfrac{h}{\sqrt{2mUe}}$

\begin{itemize}
\item 波动\quad 多\quad \unclear{概率}
\item 运动\quad 少\quad \unclear{相关}
\end{itemize}
% 原稿：“波粒二象性”下方两条箭头分别指向“波动”与“运动”；“波动”下有小字“多”及两字（疑“概率”）；“运动”下有小字“少”及“相关”

$\Delta x\,\Delta p\ge\dfrac{h}{4\pi}$\quad 概率波

电子波不能用牛顿第二定律，（无轨迹，无坐标）\ \red{（微观高\unclear{速}…} % 红字右侧被照片边缘截断

非经典力学问题\ \red{无法用牛顿运动定律}

% ---- block 192-05 (原稿第192页) kind=knowledge ----
5.\ $\Delta x\,\Delta p\ge$\unclear{$\dfrac{h}{4\pi}$}\ 不可能同时确定微观粒子的位置和动量。 % 原稿“位置和动量”被圈起；“≥”后的式子被笔画覆盖难辨，疑为h/4π；此行有一道横线贯穿


\section{氢原子光谱与玻尔模型}
\subsection{光谱与氢原子光谱}
% ---- block 004-02 (原稿第4页) kind=summary ----
\textbf{光谱}
\begin{itemize}
\item 发射光谱
  \begin{itemize}
  \item 连续谱
  \item \fbox{线状谱}
  \end{itemize}
\item \fbox{吸收光谱}
\end{itemize}
\noindent \fbox{线状谱}、\fbox{吸收光谱}$\;\to\;$原子的特征谱线 % 原稿：“线状谱”与“吸收光谱”的引线合并后，以一个箭头指向“原子的特征谱线”

% ---- block 113-02 (原稿第113页) kind=summary ----
% CHECK: 最左侧竖排三字疑为“连续谱”，其后大括号内为“吸收谱/发射谱”，层级关系按原稿位置转录
连续谱\ $\left\{\begin{tabular}{@{}l@{}}吸收谱\\发射谱\end{tabular}\right.$\ 反映内部结构光谱\ $\left\{\begin{tabular}{@{}l@{}}连续谱\\离散谱\end{tabular}\right.$

% ---- block 004-03 (原稿第4页) kind=formula ----
\[
\frac{1}{\lambda}=R\left(\frac{1}{2^2}-\frac{1}{n^2}\right)\qquad\bigl(n=3,4,5\ldots\quad \overset{\text{里德伯常量}}{R=1.10\times10^{7}\unit{m^{-1}}}\bigr)
\] % CHECK: 第二项分母手写形似“2^n”，按巴耳末公式转为 n^2；“里德伯常量”写在 R=… 上方

% ---- block 006-03 (原稿第6页) kind=formula ----
\textbf{巴尔末系}\quad $\dfrac{1}{\lambda}=R\left(\dfrac{1}{2^2}-\dfrac{1}{n^2}\right)$

% ---- block 005-04 (原稿第5页) kind=note ----
\noindent $\langle$H光谱只有4种可见光\\
\hspace*{1.5em}$6\to2$\quad $5\to2$\\
\hspace*{1.5em}$4\to2$\quad $3\to2$ $\rangle$ % 原稿：两侧为一对大括号/尖括号

\subsection{玻尔模型与氢原子能级}
% ---- block 006-04 (原稿第6页) kind=knowledge ----
\textbf{玻尔能级模型}\quad 三条假设：
\begin{itemize}
\item[①] \textsuperscript{原子}轨道/能量（能级）\textsuperscript{分立的}不连续\quad 电子在\textsuperscript{加速}高速运动但不向外辐射能量（电磁波）\quad 称为定态 % “加速”二字小字写在“高速”上方，疑为对“高速”的订正/补充
\item[②] 原子从一定态\textsuperscript{跃迁}到另一定态\quad 吸收或辐射\unclear{一定}的电磁波能量\quad 由两个能级差决定
\item[③] 原子能级$=$电子动能$+$系统电势能
\end{itemize}

（原子跃迁$=$电子轨道的跃迁） % 原稿“原子”与“跃迁”之间有一小块涂改液涂白

高轨低速大周期、大机大势\textsubscript{电势能}\ 大能量\textsubscript{能级}

% ---- block 006-05 (原稿第6页) kind=knowledge ----
\textbf{氢原子能级图}

\[
r_n=n^2r_1,\qquad E_n=\frac{1}{n^2}E_1\quad(E_1=-13.6\unit{eV})
\]

\notefig[0.35]{figures/p006_a.png}

氢原子只发出4种可见光\ $\to$\ $6\to2$\quad $5\to2$\quad $4\to2$\quad $3\to2$ % 原稿一条长箭头从“只发出4种可见光”指向页面右上方的“6→2  5→2  4→2  3→2”

\notefig[0.25]{figures/p006_b.png}

间隙越来越大 % 图 p006_b（半径示意图）内亦含此字

\notefig[0.45]{figures/p006_c.png}

间距越来越小 % 图 p006_c（能级示意图）内亦含此字

$n=1$叫基态\quad $n>1$叫激发态

$n=2$\ 第一激发态…

% ---- block 006-06 (原稿第6页) kind=formula ----
能级$\ge0$叫电离\qquad 能级间距$=$代表吸收或辐射电磁波/光子

\[
h\nu_1=h\nu_2+h\nu_3\qquad \nu_1=\nu_2+\nu_3 % CHECK: 右式等号手写形似“≠”，按物理关系应为“=”，照原样按“=”转录
\]
\[
h\frac{c}{\lambda_1}=h\frac{c}{\lambda_2}+h\frac{c}{\lambda_3}\qquad \frac{1}{\lambda_1}=\frac{1}{\lambda_2}+\frac{1}{\lambda_3}
\]

% ---- block 190-07 (原稿第190页) kind=formula ----
\[
\frac{1}{\lambda_{31}}=\frac{1}{\lambda_{32}}+\frac{1}{\lambda_{21}}
\]

% ---- block 002-06 (原稿第2页) kind=note ----
\textbf{原子的核式结构}

库仑定律、牛顿第二定律、圆周运动公式、功能关系、能量守恒定律、

% ---- block 188-02 (原稿第188页) kind=note ----
2、电子跃迁到低能级 $E_k\uparrow$ $E_{\text{电}}\downarrow$ $E_{\text{原子}}\downarrow$ 释放光子．% 注：“E原子↓”写在“E电↓”右下方、“释放光子”下方，带向下箭头

% ---- block 189-10 (原稿第189页) kind=note ----
% 原稿此块被划去：一条横线穿过“E=hc/λ=hν”并延伸到“低轨”二字，“E小”未被划
\struck{$E=\dfrac{hc}{\lambda}=h\nu$}\quad 低轨 $E$ 小

\subsection{能级跃迁与光谱线}
% ---- block 192-03 (原稿第192页) kind=knowledge ----
3.\ 高能级$\to$低能级\ 辐射发出光子，低能级$\to$高能级\ 吸收光子。

% ---- block 002-07 (原稿第2页) kind=knowledge ----
\textbf{玻尔理论的理解与应用}

\textbf{两类能级跃迁.} % 页面左侧标注，其右为大括号，含下面两项
\begin{itemize}
\item 自发跃迁（高$\to$低）
\item 受激跃迁（低$\to$高）\ $\xrightarrow{\ \text{吸}E\ }$
  \begin{itemize}
  \item 光照（吸收光子）：必须恰好 $h\nu=E_m-E_n$
  \item 碰撞加热.\quad $E_{\text{外}}\ge E_m-E_n$
  \item 吸收大于电离能的光子：原子发生电离 $E_\infty$.
  \end{itemize}
\end{itemize}

% 页面右侧另有一个大括号，含下面三项
\begin{itemize}
\item 电离态
\item $n=\infty$\quad $E=0$. % 原稿“E”前有一个极小的“2”形笔画，疑为逗号/污迹，未转录
\item 电离能\quad 原子从基态或某一激发态跃迁到电离态所需要吸收的最小能量.
\end{itemize}

% ---- block 007-01 (原稿第7页) kind=knowledge ----
\textbf{电子轨道跃迁}

一群\,/\,一个\quad $n$能级氢原子向低能级跃迁时，辐射光子的 能量/频率/波长/光谱 有\quad $\left\{\begin{array}{l}C_n^2\ \text{种}\\ n-1\ \text{种}\end{array}\right.$ % 原稿左侧“一群”(上)“一个”(下)两行共用一个大括号；右侧大括号上行“$C_n^2$种”、下行“$n-1$种”，分别与“一群”“一个”上下对应

高$\to$低\qquad 低$\to$高（吸收） % 原稿此两处写在上一行下方，“高→低”在左，“低→高(吸收)”偏右

光照跃迁\quad 不多不少\quad $h\nu=|E_{\text{末}}-E_{\text{初}}|$

轰击跃迁\quad $E_k>|E_{\text{末}}-E_{\text{初}}|$（不相等）\quad 只多不少

光照电离\quad 只多不少\quad $h\nu\ge|E_{\text{初}}|=|0-E_{\text{初}}|$（可以相等）

轰击电离\quad 只多不少\quad $E_k>|0-E_{\text{初}}|=|E_{\text{初}}|$
% 原稿左侧有一个大括号把“轰击跃迁、光照电离、轰击电离”三行括在一起（“光照跃迁”在括号之外）

% ---- block 003-01 (原稿第3页) kind=knowledge ----
\textbf{氢原子辐射光谱线的$\langle$数量$\rangle$} % 原稿：“数量”二字被一对大尖括号圈起(强调)
\begin{itemize}
\item[] 一个氢原子\quad $n-1$ 条
\item[] 一群氢原子\quad $C_n^2$ 条\quad $\dfrac{n(n-1)}{2}$
\end{itemize}

% ---- block 191-08 (原稿第191页) kind=knowledge ----
8.\ 氢原子从高能级向低能级跃迁不产生$\gamma$射线，（$\gamma$射线是原子核能级跃迁时产生的）

\hspace{1.5em}从$n=4$向低能级跃迁时辐射出6种频率的光，只有2种可见光。

\hspace{1.5em}红外线\ 可见光\ 紫外线是原子外层电子受激发产生的

\hspace{1.5em}$X$射线是原子内层电子受激发产生的。$\gamma$射线是原子核的能级跃迁产生的

% ---- block 188-13 (原稿第188页) kind=note ----
10、一群氢原子有 $\dfrac{n(n-1)}{2}$ 种光子\quad 1 个氢原子有 $n-1$ 种光子．（光电效应仅表粒子性）．

$E_k=h\nu-W_0$\quad $p=\dfrac h\lambda$\quad $p=\dfrac Ec$\quad 故光子能量越大 动量也越大．\quad $\lambda=\dfrac hp=\dfrac{h}{\sqrt{2mE}}$．

% ---- block 189-11 (原稿第189页) kind=note ----
% 序号被页边遮挡（疑为 8）
$n=2\to n=1$ 放出光子恰发生光电效应．从 $n=3\to n=1$\quad $E_{k\text{初}\max}=E_3-E_2$

$W_0=E_2-E_1$\quad\quad $E_k=E_3-E_1-(E_2-E_1)$

% ---- block 191-04 (原稿第191页) kind=problem ----
4.
\begin{problem}
$E_n=\dfrac{E_1}{n^2}$\quad $n=2,3,4$。H从$n=2\to n=1$放出光子频率$\nu$，巴尔末系指H从高能级向$n=2$能级跃迁时释放的光子，使

使处于基态的氢原子释放出巴尔末系的光子的照射光光子最小频率？ % CHECK: 上一行末尾的“使”与本行开头的“使”重复，照原样转录
\begin{solution}
\[
\frac{E_1}{4}-E_1=h\nu
\]
\[
\frac{E_1}{9}-E_1=h\nu'\qquad \nu'=\frac{32}{27}\nu
\]
\end{solution}
\end{problem}


\section{天然放射现象与原子核的组成}
% ---- block 007-02 (原稿第7页) kind=summary ----
\textbf{原子核衰变过程}\quad 放出$\alpha$、$\beta$、$\gamma$（射线）称天然放射现象（贝克勒尔发现）\quad \textsuperscript{原子}核是复杂的

\begin{center}
\small\setlength{\tabcolsep}{4pt}
\begin{tabular}{c l c c c c c l}
\toprule
 & \textsuperscript{本质} & 质量 & 电量 & 电磁偏转 & $v$ & \makecell{电离性\\ $\alpha$\unclear{m}} & \makecell{穿透性\\ $\alpha$\unclear{$\vee$}} \\
\midrule
$\alpha$ & $\nuc{4}{2}{He}$\ $\alpha$粒子 & $4M$ & $+2e$ & 能 & 小 & \multirow{3}{*}{$\uparrow$\,强} & 纸 \\
$\beta$ & $\nuc{0}{-1}{e}$\ $\beta$粒子 & $\approx0$ & $-e$ & 能 & 中 & & 铝板 \\
$\gamma$ & 光子\ 电磁波 & $0$ & $0$ & 不能 & 大 & & $\downarrow$\,强\ \unclear{水泥墙} \\
\bottomrule
\end{tabular}
\end{center}
% 原稿：“电离性”列下有一条自下而上贯穿三行的长箭头，箭头上端写“强”；“穿透性”列有一条自上而下贯穿三行的长箭头，箭头下端写“强”，三行分别为“纸”“铝板”“水泥墙”（该行末字手写不清）；表格左侧有一条竖线把α、β、γ与右侧隔开
% CHECK: “电离性”“穿透性”两个表头后面的小字（形似“αm”“α∨”）含义不明，照原样标为不清

% ---- block 190-01 (原稿第190页) kind=knowledge ----
1.\ 阴极射线说明原子可分有内部结构

\unclear{$\gamma$}射线是原子核由高能级向低能级跃迁时发出的高频光子流。（实质\ \textit{核中中子变为质子}\ $\nuc{1}{0}{n}\to\nuc{1}{1}{H}+\nuc{0}{-1}{e}$）形成电子流。 % 原稿"核中中子变为质子"小字写在括号内公式上方 % CHECK: 句首字被改写覆盖(γ/β)，括号内为β衰变实质

% ---- block 192-02 (原稿第192页) kind=knowledge ----
2.\ $\gamma$射线与$\nuc{1}{0}{n}$不会发生电磁偏转。（因$\gamma$与$\nuc{1}{0}{n}$不带电）

% ---- block 191-02 (原稿第191页) kind=knowledge ----
2.\ 用电磁场来加速减速来\unclear{束缚}中子\ $\times$\qquad 中子不带电不受电磁场作用 % 原稿“来”后一字难辨(缚/束缚?)

% ---- block 192-06 (原稿第192页) kind=knowledge ----
\underline{原子核不存在独立的电子}（仅有质子中子）


\section{原子核的衰变与半衰期}
% ---- block 002-08 (原稿第2页) kind=summary ----
\textbf{原子核的衰变及半衰期}

\subsection{衰变的类型与衰变次数}
% ---- block 007-04 (原稿第7页) kind=knowledge ----
\textbf{衰变}
\begin{itemize}
\item 一变多：天然衰变，自发的
\item 83号以上元素都可以衰变，83号以下元素部分可衰变\quad $\nuc{1}{1}{H}$\quad $\nuc{17}{8}{O}\ \checkmark$\quad $\nuc{16}{8}{O}\ \times$ % “$^{17}_{8}$O”的上标17手写不太清（也可能是18）
\end{itemize}

\textbf{$\alpha$衰变}\quad 本质：核子的重新组合\qquad 都逆时针转

\hspace{2em}一变多\textsuperscript{二}，生成物中有$\alpha$粒子（氦核$\nuc{4}{2}{He}$）\quad 退2少4\ 进1保质\quad $R=\dfrac{p}{qB}$\ 外切 % “多”字上方有一个小“=”（疑为“二”）

\textbf{$\beta$衰变}\quad 本质\ $\nuc{1}{0}{n}\to\nuc{1}{1}{H}+\nuc{0}{-1}{e}$（分裂）\qquad $R\propto\dfrac{1}{q}$ % “$R\propto 1/q$”原稿写在α、β两行右侧之间

\hspace{2em}一变多，生成物中有$\beta$粒子（电子$\nuc{0}{-1}{e}$）\quad 电子顺新核逆时针\quad $R=\dfrac{p}{qB}$\ 内切

\textbf{$\gamma$衰变}\quad $\gamma$射线：电磁波（光子）\textsuperscript{仅能}伴随$\alpha$衰变和$\beta$衰变产生

% ---- block 003-02 (原稿第3页) kind=knowledge ----
由质量数守恒定$\alpha$衰变次数\quad 再由电荷数守恒确定$\langle$$\beta$衰变次数$\rangle$ % 原稿：“β衰变次数”被一对大尖括号圈起

\noindent $\alpha$衰变\quad $\nuc{A}{Z}{X}\to\nuc{A-4}{Z-2}{Y}+\nuc{4}{2}{He}$\qquad $\beta$衰变\quad $\nuc{A}{Z}{X}\to\nuc{A}{Z+1}{Y}+\nuc{0}{-1}{e}$

\noindent $2\,\nuc{1}{1}{H}+2\,\nuc{1}{0}{n}\to\nuc{4}{2}{He}$\qquad $\nuc{1}{0}{n}\to\nuc{1}{1}{H}+\nuc{0}{-1}{e}$ % CHECK: 第一式中“H”的上下标位置不清，按质子 1/1 H 转录

% ---- block 188-11 (原稿第188页) kind=formula ----
$\beta$ 衰变\quad $\nuc{1}{0}{n}\to\nuc{1}{1}{H}+\nuc{0}{-1}{e}$

% ---- block 190-04 (原稿第190页) kind=knowledge ----
% 原稿：整句下划线；“α、β衰变”、“一定”、“有γ放出”分别被圈起
4.\ \unclear{$\alpha$}、$\beta$衰变一定有$\gamma$放出

{\small （\unclear{伴$\gamma$射线}）} % 原稿此小字写在“有γ放出”正下方，字迹很小

\subsection{衰变与动量、磁场中的径迹}
% ---- block 004-01 (原稿第4页) kind=method ----
\textbf{$\langle$衰变+动量$\rangle$} % 原稿：标题被一对大尖括号圈起

\noindent $\alpha$衰变.\quad $\nuc{A}{Z}{X}\to\nuc{A-4}{Z-2}{Y}+\nuc{4}{2}{He}$\quad \fbox{退2\unclear{少}4} % 圈内文字“退2少4”中“少”字不清(也可能是“以”)

两圆外切，$\alpha$粒子半径较大

\notefig[0.25]{figures/p004_a.png}
\begin{gather*}
m_1v_1-m_2v_2=0\\
E=\frac{p^2}{2m_1}+\frac{p^2}{2m_2}\\
r=\frac{p}{qB}
\end{gather*}
\[
I=\frac{q}{T}=\frac{q\,qB}{2\pi m}=\frac{q^2B}{2\pi m}
\]

\noindent $\beta$衰变\quad $\nuc{A}{Z}{X}\to\nuc{A}{Z+1}{Y}+\nuc{0}{-1}{e}$\quad \fbox{进1\unclear{位}} % 圈内文字“进1位”中“位”字不清

两圆内切，$\beta$粒子半径较大.

\notefig[0.25]{figures/p004_b.png}

% ---- block 089-05 (原稿第89页) kind=method ----
\noindent $\alpha$ 衰变\quad 外切
\notefig[0.25]{figures/p089_d.png}
\[ R_\alpha=\frac{m_\alpha V_\alpha}{q_\alpha B}\qquad R_{\text{新}}=\frac{m_{\text{新}}V_{\text{新}}}{q_{\text{新}}B} \]
\begin{keybox}
$R\propto\dfrac1q$\quad 故 $R_\alpha$、$R_\beta$ 大。
\end{keybox}
\noindent $\beta$ 衰变\quad 内切
\notefig[0.25]{figures/p089_e.png}

\subsection{半衰期}
% ---- block 003-03 (原稿第3页) kind=formula ----
\[
m_{\text{余}}=m_{\text{原}}\left(\frac12\right)^{\frac{t}{\tau}}
\] % CHECK: 指数分母(半衰期)符号不清，似“τ”或“T”

% ---- block 008-03 (原稿第8页) kind=knowledge ----
\textbf{半衰期}\quad 未衰变物质质量\ $m=M_0\left(\dfrac{1}{2}\right)^n$\qquad $n$个半衰期

\hspace{4.5em}已衰变物质质量\ $M_0-m$\qquad\qquad 说克不说个，是大量

\unclear{衡}量衰变快慢的物理量\qquad 16个\unclear{氡}经1个半衰期，还剩8个\ $\times$

只由元素决定\qquad 16g\unclear{氡}经1个半衰期，还剩$\underset{\approx16\unit{g}}{8\unit{g}}$\ $\times$\quad 还剩$8\unit{g}$未衰变
% 原稿“还剩8g”下方有小字“≈16g”；“氡”字（疑）原稿写成圈起来的“80”状字形；左侧两行“（衡）量衰变快慢的物理量”“只由元素决定”位于下方两行的左边

% ---- block 188-09 (原稿第188页) kind=note ----
8、核聚变不可控．\quad 衰变公式 $\dfrac{m_{\text{余}}}{m}=\left(\dfrac12\right)^{\frac nT}$% CHECK: 左侧分数字迹潦草，疑为 m余/m；指数疑为 n/T
\quad（$n$ 为时\unclear{间}，$T$ 为半衰期）．经 $2T$ 4 个 \unclear{U}238 变为 2 个% 原稿此句上有红色斜线划过

\red{半衰期是本征属性，只与原子核有关，与外界无关}\quad \red{对大量原子核的统计规律}

（质量 \unclear{1个} 减 1 半 $\times$，不能对少量物质说，只能广泛地说）


\section{核反应方程}
% ---- block 002-09 (原稿第2页) kind=summary ----
\textbf{核反应方程与核能的计算.}

% ---- block 003-04 (原稿第3页) kind=summary ----
\textbf{核反应四种类型}\quad 可控性 % 原稿：“可控性”写在右侧，作为后面“自发”“人工控制”等的列标题
\begin{itemize}
\item 衰变
  \begin{itemize}
  \item $\alpha$衰变\quad 自发\quad $\nuc{238}{92}{U}\to\nuc{234}{90}{Th}+\nuc{4}{2}{He}$
  \item $\beta$衰变\quad 自发\quad $\nuc{234}{90}{Th}\to\nuc{234}{91}{Pa}+\nuc{0}{-1}{e}$
  \end{itemize}
\item $\langle$人工转变$\rangle$\quad 人工控制
  \begin{itemize}
  \item $\nuc{14}{7}{N}+\nuc{4}{2}{He}\to\nuc{17}{8}{O}+\nuc{1}{1}{H}$\quad 卢瑟福发现质子
  \item $\nuc{4}{2}{He}+\nuc{9}{4}{Be}\to\nuc{12}{6}{C}+\nuc{1}{0}{n}$\quad 查德威克发现中子
  \item $\nuc{27}{13}{Al}+\nuc{4}{2}{He}\to\nuc{30}{15}{P}+\nuc{1}{0}{n}$\quad 居里夫妇发现放射性同位素
  \item $\nuc{30}{15}{P}\to\nuc{30}{14}{Si}+\nuc{0}{+1}{e}$\quad 同时发现正电子
  \end{itemize}
\item $\langle$重核裂变$\rangle$\quad 易人工控制
  \begin{itemize}
  \item $\nuc{235}{92}{U}+\nuc{1}{0}{n}\to\nuc{144}{56}{Ba}+\nuc{89}{36}{Kr}+3\,\nuc{1}{0}{n}$
  \item $\nuc{235}{92}{U}+\nuc{1}{0}{n}\to\nuc{136}{54}{Xe}+\nuc{90}{38}{Sr}+10\,\nuc{1}{0}{n}$
  \end{itemize}
\item $\langle$轻核聚变$\rangle$\quad 很难控制\quad $\nuc{2}{1}{H}+\nuc{3}{1}{H}\to\nuc{4}{2}{He}+\nuc{1}{0}{n}$
\end{itemize}

% ---- block 007-03 (原稿第7页) kind=knowledge ----
质量数$\to\nuc{235}{92}{U}$\ $\leftarrow$电荷量 % 原稿：“质量数”的箭头指向上标235，“电荷量”的箭头由左下方指向下标92

\begin{center}
{\renewcommand{\arraystretch}{1.6}
\begin{tabular}{|c|c|c|}
\hline
$\nuc{1}{1}{H}$ & $\nuc{2}{1}{H}$ & $\nuc{3}{1}{H}$ \\ \hline % 原稿第三格“3H”的下标1不明显
$\nuc{1}{0}{n}$ & $\nuc{0}{-1}{e}$ & $\nuc{4}{2}{He}$ \\ \hline
\end{tabular}}
\end{center}

\textbf{核反应}\quad 三变\ 四恒\ 一\unclear{奇葩}

守恒：质量数、电荷数、能量、动量

不守恒：质子数、中子数、电子数（$\beta$衰变）

质量：质量亏损$\checkmark$\ （静态质量）\quad 质量守恒$\checkmark$\ （相对论质量）\quad $E=mc^2$ % 原稿两处括号写在对应词的下一行

$M_{\text{反}}\neq M_{\text{生}}$

% ---- block 191-06 (原稿第191页) kind=knowledge ----
6.\ 核反应遵循\underline{电荷守恒}\ \underline{质量数守恒}\ 但\underline{质量不守恒}

% ---- block 007-05 (原稿第7页) kind=knowledge ----
\textbf{人工转变}
\begin{itemize}
\item 二变二。非自发\textsuperscript{核}反应 % 原稿“二变=.”，末字“二”写作“=”
\item 反应物中有$\alpha$粒子\ $\nuc{4}{2}{He}$
\end{itemize}

\[
\nuc{4}{2}{He}+\nuc{14}{7}{N}\to\nuc{17}{8}{O}+\nuc{1}{1}{H}\qquad\text{卢瑟福发现质子}
\]
\[
\nuc{4}{2}{He}+\nuc{9}{4}{Be}\to\nuc{12}{6}{C}+\nuc{1}{0}{n}\qquad\text{查德威克发现中子}
\]

% ---- block 191-12 (原稿第191页) kind=knowledge ----
卢瑟福发现质子\quad 查德威克发现中子\qquad 慢中子与\unclear{${}^{6}$Li}反应\ \underline{是人工转变}\ \underline{不是裂变反应} % 原稿“Li”左上角的数字形似“5”，按6猜


\section{核裂变与核聚变}
% ---- block 008-01 (原稿第8页) kind=knowledge ----
\textbf{重核裂变}（链式反应）\quad 1中子$\to$多于2个中子\qquad 条件：体积足够大到临界体积

\[
\nuc{1}{0}{n}+\nuc{238}{92}{U}\to\mathrm{Sr}+\mathrm{Ba}+3\,\nuc{1}{0}{n}\qquad\text{原子弹、核电站}
\]
% CHECK: $\mathrm{U}$ 的质量数手写形似238（常见裂变用235）；第一个产物符号形似“Sr”（亦可能是Kr），均照原样转录

% ---- block 008-02 (原稿第8页) kind=knowledge ----
\textbf{轻核聚变}\quad 两小生一大，中子是附加\unclear{品}\qquad 条件：几百万度高温，热核反应

\[
\nuc{2}{1}{H}+\nuc{3}{1}{H}\to\nuc{4}{2}{He}+\nuc{1}{0}{n}\qquad\text{（用原子弹引爆）}\quad\text{也叫轻核骤变}
\]
% CHECK: “骤变”疑为“聚变”（原稿字形为“骤”）

% ---- block 189-04 (原稿第189页) kind=note ----
% 原稿对“热核反应”“核聚变”“除氢弹外全为核裂变”“镉棒控制核反应速度”“临界质量才裂变”有下划线，“铀块”与“大于”各被圈出；“放能”被红叉划去
4、热核反应是核聚变．除氢弹外全为核裂变\quad 镉棒控制核反应速度．（\unclear{铀块}大于临界质量才裂变）

一切核反应都\struck{放能}\quad \red{只有有质量亏损的反应才放能}

% ---- block 229-03 (原稿第229页) kind=note ----
裂变聚变都放能，

\unclear{核}铀裂变释放 快中子不易捕获 衰变自发 不需要中子 % CHECK: "铀"字上方有一个小字（插入），辨认不清


\section{核力、结合能与质量亏损}
\subsection{核力与结合能}
% ---- block 008-04 (原稿第8页) kind=knowledge ----
\textbf{核力}\quad 核子与核子\ $\downarrow$\ 强力（\unclear{交换介子}）\ $\downarrow$ % 原稿竖排：“核子与核子”下接箭头指向“强力（…）”，其下再有一箭头指向“结合能”；括号内字被箭头竖线穿过，难辨

\begin{itemize}
\item 核子\fbox{分开}强力做负功，\fbox{吸收}能量，$m$增大
\item 核子\fbox{结合}强力做正功，\fbox{放出}能量，$m$减少
\end{itemize}
% 以上两行原稿共用一个大括号（括号左侧即“强力（…）”）；加框的词在原稿中是手绘圆圈圈出

\notefig[0.3]{figures/p008_a.png}

结合能$=$比结合能$\times$核子数

（键能）\quad $\Delta E=\Delta m\cdot c^2$

核子数越多，结合能越大\quad Fe % “Fe”小字写在这一行末尾右侧上方（对应上图曲线峰值处的Fe）

比结合能（结合能/核子数量）：比结合能越大，原子核越稳定

% ---- block 190-06 (原稿第190页) kind=knowledge ----
新核核子平均质量小，比结合能增大。

% ---- block 285-05 (原稿第285页) kind=formula ----
% 第一个核素的原子序数左侧有一个小圈状记号
\[ \nuc{A}{Z}{X}\to\nuc{A'}{Z'}{X'}+\nuc{A''}{Z''}{Y'}\qquad \text{放能}=A'E_2+A''E_3-AE_1, \]
\[ AE_1<A'E_2+A''E_3 \]

% ---- block 188-04 (原稿第188页) kind=note ----
3、稳定性 = 比结合能（与结合能无关）

% ---- block 188-08 (原稿第188页) kind=note ----
7、生成物比结合能 大于 反应物（不是化学反应）．对放能核反应 生成物结合能也更大

生成 He 为 $\alpha$ 衰变，生成 $\mathrm{e}^-$ 为 $\beta$ 衰变，生成 $\gamma$ 为 $\gamma$ 衰变\quad 中子 $\nuc{1}{0}{n}$\quad $\nuc{0}{-1}{e}$

\subsection{质量亏损、质能方程与核能计算}
% ---- block 192-04 (原稿第192页) kind=knowledge ----
4.\ \underline{原子核静质量}小于\underline{全部核子独存时的总静质量}。 % 原稿“小于”被圈起

% ---- block 193-04 (原稿第193页) kind=formula ----
质能方程\quad $E=\Delta mc^2$\qquad $E=h\nu=\dfrac{hc}{\lambda}$\qquad $\lambda=\dfrac{hc}{E}$

% ---- block 003-05 (原稿第3页) kind=formula ----
\textbf{$\langle$核能计算$\rangle$} % 页面左侧标注
\begin{itemize}
\item[$\downarrow$] $\Delta E$与$\Delta m$成正比 但不能说相互转化
\item[] $\Delta E=\Delta mc^2$\quad $\Delta m$单位 $\mathrm{kg}$\quad $c$单位 $\mathrm{m/s}$\quad $\Delta E$单位 $\mathrm{J}$.
\item[] $\Delta E=\Delta m\times931.5\unit{MeV}$\quad $\Delta m$单位为 $\mathrm{u}$，$\Delta E$单位为 $\mathrm{MeV}$
\item[] 原子核的结合能 = 核子比结合能 $\times$ 核子数.
\end{itemize}

% ---- block 008-05 (原稿第8页) kind=knowledge ----
\[
\nuc{2}{1}{H}+\nuc{3}{1}{H}\to\nuc{4}{2}{He}+\nuc{1}{0}{n}+17.6\unit{MeV}
\]

\notefig[0.27]{figures/p008_b.png}

$1.09\unit{MeV}$\quad $2.78\unit{MeV}$ % 原稿写在上式两个反应物的下方，并带两条弧形箭头（见图 p008_b；图内已含此二数）

\[
E_{\text{前}}+\underset{17.6\unit{MeV}}{E_{\text{放}}}=E_{\text{后}}\qquad Q=\Delta mc^2
\]

若\unclear{$\alpha$}以$Q$轰击$\nuc{14}{7}{N}$则不会反应，还需提供动能 % “若”后第一个字（形似“X”）疑为α

生成物总结合能$>$反应物总结合能\ \textsuperscript{\unclear{光子有动质量}}

新核和$\alpha$粒子动量不等大反向\quad 光子也有动量\quad 释放能量给了$\gamma$光子和新核的动能

不能说质量就是能量，不能表示质量和能量间转化关系\quad 可说正比

% ---- block 189-01 (原稿第189页) kind=note ----
1、1个\unclear{核反应}释放\unclear{的} $E_0$ 能量\quad $\mathrm{MeV}=E_0\times1.6\times10^{-13}\unit{J}$


\section{实验}
% ---- block 005-07 (原稿第5页) kind=summary ----
\noindent \textbf{$\langle$三个电路$\rangle$}

\noindent 阻止\ （\unclear{光向}正极板）\ 光强不变，仅提高照射光频率，光子数小，电流$\downarrow$ % “光向”两字被覆写，不清；该行位于电路②③标题上方

% ---- block 005-08 (原稿第5页) kind=diagram ----
\noindent \red{①} 光电效应发生电路

\notefig[0.3]{figures/p005_b.png}
\noindent 瞬时性（不要时间） % CHECK: 首字手写形似“目急”，按语义转为“瞬”

% ---- block 005-10 (原稿第5页) kind=diagram ----
\noindent \red{②} 遏止电压 $U_c$ 测定电路

\notefig[0.35]{figures/p005_c.png}
\noindent 负极接收\textsuperscript{极}$\to$遏止\\ % 原稿：“极”字小写在“收”上方，为插入
看到遏止电极压 $U_c$ % CHECK: 疑为“遏止电压”，“极”字原稿可见，照录
\[
\begin{array}{l}
-eU_c=0-E_k\\
\hookrightarrow\ eU_c=E_k=h\nu-W_0=h\dfrac{c}{\lambda}-W_0
\end{array}
\]

% ---- block 005-11 (原稿第5页) kind=diagram ----
\noindent \red{③} \underline{饱和光电流测定电路}\ 光向负极板

\notefig[0.3]{figures/p005_d.png}
\noindent 原因：单位时间 $K$ 发出光电子数一定，达到某个电压后，光电子全部\unclear{被}阳极吸收，电流不再改变 % 原稿：“原因”二字写在图中阳极板右上角；“阳极”为行间小字插入；图中“hν”处引出一条弧线箭头指向“单位时间K”

\noindent 正极接收板$\to$促进\\
饱和光电流 只和 $n$ 有关\ （光子数）

\noindent 光强 $P=nh\nu$\\
$n$ 单位时间光子数\ （电子数）

